% Baccalauréat malgache, série C, 2016 : énoncé.
\TCDBacAttribution{C}{2016}
\providecommand{\cm}{~\mathrm{cm}}
\providecommand{\D}{(\mathscr{D})}
\providecommand{\Courbe}{(\mathscr{C})}
\providecommand{\CourbeP}{(\mathscr{C}')}
\providecommand{\Plan}{(\mathscr{P})}
\providecommand{\mes}[2]{\mathrm{mes}\left( \widehat{ \overrightarrow{#1} ; \overrightarrow{#2} } \right)}
\providecommand{\Angl}[2]{\left( \overrightarrow{#1} ; \overrightarrow{#2} \right)}
%%%%%%%%%%%%%%%%%%%%%%%%%%%%%%%%%%%%%%%%%%%%%%%%%%%%%%%%%%%%%%%%%%%%%%%%%%%%%%%%%%%%%%%%%%%%%%%%%%%%%%%%%%%%%%%%%%%%%%%%%%%%%%%%
\noindent \underline{\textsc{Exercice}} (4 points)% \newline
      \begin{description}
            \item[I.\underline{Arithm\'etique :}]
                  \begin{enumerate}
                        \item On consid\`ere l'\'equation \`a deux inconnues $(x,y)$ dans $\mathbb{Z}\times\mathbb{Z}$ suivante :
                              \begin{equation}
                                    \TCDtag{E}
                                          3x - 8y = 1 \;.
                                    \label{eq: Exo Arithm Bac C 2016}
                              \end{equation}
                              \begin{enumerate}
                                    \item V\'erifier que $(3;1)$ est une solution de \eqref{eq: Exo Arithm Bac C 2016}.
	                                    \item R\'esoudre l'\'equation \eqref{eq: Exo Arithm Bac C 2016}.
                              \end{enumerate}
                        \item $a$, $b$, $c$ sont des entiers natures non nuls.
                              Un entier nature non nul $A$ s'\'ecrit $\overline{(b0a)}_{5}$ dans le syst\`eme de num\'eration
                              \`a base cinq et $\overline{(abc)}_{7}$ dans le syst\`eme de num\'eration \`a base sept.
                              D\'emontrer que $c = 6$ puis d\'eterminer $a$ et $b$.
                  \end{enumerate}
            \item[II.\underline{Probabilit\'e :}]
                  Le code confidentiel d'une carte bancaire est un nombre entier de quatre chiffres non nuls.
                  \begin{enumerate}
                        \item Combien y a-t-il de codes possibles\,?
                        \item Le code d'une carte est choisi au hasard.
                              Calculer les probabilit\'es de chacun des \'ev\`enements suivants :
                              \begin{itemize}
                                    \item $A$ : << Le code est un nombre pair >>.
                                    \item $B$ : << Le code n'est compos\'e que de chiffres pairs >>.
                                    \item $C$ : << Le code contient une fois et une seule le chiffre 2 >>.
                                    \item $D$ : << Le code est \'ecrit avec des chiffres distincts >>.
                              \end{itemize}
                  \end{enumerate}
      \end{description}
%%%%%%%%%%%%%%%%%%%%%%%%%%%%%%%%%%%%%%%%%%%%%%%%%%%%%%%%%%%%%%%%%%%%%%%%%%%%%%%%%%%%%%%%%%%%%%%%%%%%%%%%%%%%%%%%%%%%%%%%%%%%%%%%

%%%%%%%%%%%%%%%%%%%%%%%%%%%%%%%%%%%%%%%%%%%%%%%%%%%%%%%%%%%%%%%%%%%%%%%%%%%%%%%%%%%%%%%%%%%%%%%%%%%%%%%%%%%%%%%%%%%%%%%%%%%%%%%%
\noindent \underline{\textsc{Probleme I}} (7 points) \newline
      Les deux parties A et B sont ind\'ependantes.
      On note $S_{(\Delta)}$ la sym\'etrie orthogonale d'axe $(\Delta)$,
      par $t_{\vec{V}}$ la translation de vecteur $\overrightarrow{V}$.
      
      \underline{\textbf{Partie A :} M\'ethode g\'eom\'etrique} \newline
            \begin{minipage}{0.45\actualtextwidth}
                  Dans le plan orient\'e $\Plan$, on donne un carr\'e direct $ABCD$ de centre $I$, 
                  avec $\Angl{AB}{AD} = \frac{\pi}{2}$.
                  On note par $E$ le milieu de $[AB]$, $F$ le milieu de $[AD]$ et par $J$ le milieu de $[CD]$.            
            \end{minipage}
            %
            \begin{minipage}{0.45\actualtextwidth}
            \begin{flushright}
                  \begin{tikzpicture}[line cap=round,line join=round,>=triangle 45,x=.5cm,y=.5cm]
                        \draw (-2,0)-- (2,0);
                        \draw (2,0)-- (2,4);
                        \draw (2,4)-- (-2,4);
                        \draw (-2,4)-- (-2,0);
                        \draw (-2,4)-- (2,0);
                        \draw (2,4)-- (-2,0);
                        \begin{scriptsize}
                        \fill [color=black] (-2,0) circle (1.5pt) node[below left] {$A$};
                        \fill [color=black] (2,0) circle (1.5pt)node[below right] {$B$};
                        \fill [color=black] (2,4) circle (1.5pt) node[above right] {$C$};
                        \fill [color=black] (-2,4) circle (1.5pt) node[above left] {$D$};
                        \fill [color=black] (0,0) circle (1.5pt) node[below] {$E$};
                        \fill [color=black] (-2,2) circle (1.5pt) node[left] {$F$};
                        \fill [color=black] (0,4) circle (1.5pt) node[above] {$J$};
                        \fill [color=black] (0,2) circle (1.5pt) node[above] {$I$};
                        \end{scriptsize}
                  \end{tikzpicture}
            \end{flushright}
            \end{minipage}
            %
            \begin{enumerate}
                  \item Soit $G$ le barycentre du syst\`eme de points pond\'er\'es
                        $\{ (A;-1);(B;1);(D;1) \}$.
                        \begin{enumerate}
                              \item D\'emontrer que $G$ appartient \`a la droite $(AI)$.
                              \item D\'eterminer et construire $G$.
                              \item Soit $\Courbe$ l'ensemble des points $M$ du plan v\'erifiant : 
                                    $2 MI^2 - MA^2 = 0$.
                                    D\'emontrer que $B \in \Courbe$.
                                    D\'eterminer et construire $\Courbe$.
                        \end{enumerate}
                  \item Soit $\sigma$ la sym\'etrie centrale de centre $I$ et $f = \sigma \circ S_{(EF)}$.
                        En d\'ecomposant $\sigma$ en produit de deux sym\'etries orthogonales convenables,
                        d\'emontrer que $f = S_{(AC)} \circ t_{\frac{1}{2}\overrightarrow{AC}}$.
            \end{enumerate}

      \underline{\textbf{Partie B :} Utilisation de nombres complexes} \newline
            $ABCD$ est toujours le carr\'e de la partie A, de centre $I$; $\Angl{AB}{AD} = \frac{\pi}{2}$;
            le milieu de $[CD]$ est $J$.
            Soit $S$ une similitude directe qui transforme $A$ en $I$ et $B$ en $J$.
            $(\Gamma_1)$ est un cercle de diam\`etre $[AI]$ et $(\Gamma_2)$ est un cercle de diam\`etre $[BJ]$.
            On donne le rep\`ere orthonorm\'e direct $(A ; \overrightarrow{AB} , \overrightarrow{AD})$.
            \begin{enumerate}
                  \item Donner les affixes des points $A$, $B$, $C$, $D$, $I$ et $J$.
                  \item Ecrire l'expression complexe de $S$.
                        En d\'eduire ses \'el\'ements caract\'eristiques.
                  \item \begin{enumerate}
                              \item Pr\'eciser les centres et les rayons des cercles $(\Gamma_1)$ et $(\Gamma_2)$.
                              \item D\'emontrer que le centre $\Omega$ de $S$ est un point d'intersection de 
                                    $(\Gamma_1)$ et $(\Gamma_2)$.
                        \end{enumerate}
                  \item \begin{enumerate}
                              \item D\'eterminer le point $C'$, image de $C$ par $S$ puis le point $K$, image de $I$ par $S$.
                              \item D\'emontrer que les points $A$, $\Omega$, $K$ sont align\'es.
                        \end{enumerate}
            \end{enumerate}
%%%%%%%%%%%%%%%%%%%%%%%%%%%%%%%%%%%%%%%%%%%%%%%%%%%%%%%%%%%%%%%%%%%%%%%%%%%%%%%%%%%%%%%%%%%%%%%%%%%%%%%%%%%%%%%%%%%%%%%%%%%%%%%%

%%%%%%%%%%%%%%%%%%%%%%%%%%%%%%%%%%%%%%%%%%%%%%%%%%%%%%%%%%%%%%%%%%%%%%%%%%%%%%%%%%%%%%%%%%%%%%%%%%%%%%%%%%%%%%%%%%%%%%%%%%%%%%%%
\noindent \underline{\textsc{Probleme II}} (9 points) \newline
      Soit $f$ une fonction num\'erique d\'efinie sur $\mathbb{R}$ par :
      \[
      f(x) =
            \begin{multiform}
                  (x-1)\ln(1-x) \text{ si } x < 0 \\ 
                  \expo{\frac{x}{2}}-x-1 \text{ si } x \geq 0 \;.
            \end{multiform}
      \]
      On d\'esigne par $\Courbe$ la courbe repr\'esentative de $f$ dans un rep\`ere orthonorm\'e direct d'unit\'e $2\cm$.
      \begin{description}
            \item[A. ]
                  \begin{enumerate}
                        \item Etudier la continuit\'e et la d\'erivabilit\'e de $f$ en $x_0=0$.
                        \item \begin{enumerate}
                                    \item Calculer les limites de $f$ en $-\infty$ et en $+\infty$.
                                    \item Etudier les variations de $f$.
                                    \item D\'emontrer que l'\'equation $f(x) = 0$ admet une solution unique $\alpha\in[2;3]$.
                              \end{enumerate}
                        \item Tracer $\Courbe$.
                              (On pr\'ecisera les branches infinies et les demi-tangentes \`a l'origine du rep\`ere)
                  \end{enumerate}
            \item[B. ]
                  On se propose de trouver une valeur approch\'ee de la 
                  solution $\alpha$ de l'\'equation $f(x) = 0$ de la partie A.
                  On consid\`ere une fonction $g$ d\'efinie sur $]-1;+\infty[$ par : $g(x) = 2 \ln (x+1)$.
                  \begin{enumerate}
                        \item Etudier les variations de $g$.
                        \item D\'emontrer que $\alpha$ est une solution de l'\'equation $g(x) = x$.
                        \item D\'emontrer que pour tout $x \in [2;+\infty[$, $g(x) \in [2;+\infty[$.
                        \item On d\'efinit une suite $(U_n)_{n\in\mathbb{N}}$ par : 
                              $
                                    \begin{multiform}
                                          U_0 = 3 \\
                                          U_{n+1} = g(U_n) \text{ pour tout } n\in\mathbb{N} \;.
                                    \end{multiform}
                              $
                        \begin{enumerate}
                              \item D\'emontrer, par r\'ecurrence, que pour tout $n\in\mathbb{N}$, $U_n \geq 2$.
                              \item Prouver que pour tout $x \in [2;+\infty[$, $|g'(x)| \leq \frac{2}{3}$.
                              \item D\'emontrer, en utilisant les in\'egalit\'es des accroissements finis, que pour tout
                                    $n\in\mathbb{N}$, 
                                    $\displaystyle{
                                          \left| U_{n+1} - \alpha \right| \leq \frac{2}{3} \left| U_n - \alpha \right| \;.
                                    }$
                                    En d\'eduire que pour tout $n\in\mathbb{N}$ : $|U_n-\alpha| \leq \left(\frac{2}{3}\right)^n$.
                              \item D\'emontrer que la suite $(U_n)$ converge vers $\alpha$.
                                    D\'eterminer le plus entier naturel $n_0$ tel que $U_{n_0}$ soit une valeur 
                                    approch\'ee de $\alpha$ \`a $10^{-1}$ pr\`es.
                        \end{enumerate}
                  \end{enumerate}
            \item[C. ]
                  Soit l'\'equation diff\'erentielle 
                        \begin{equation}
                              \TCDtag{E}
                                    2y' - y = x - 1 \;.
                              \label{eq: Prob II partie C Bac C 2016}
                        \end{equation}
                  \begin{enumerate}
                        \item D\'emontrer que $f(x) = \expo{\frac{x}{2}} - x - 1$ 
                              est solution de \eqref{eq: Prob II partie C Bac C 2016}.
                        \item En d\'eduire la solution g\'en\'erale de l'\'equation \eqref{eq: Prob II partie C Bac C 2016}.
                  \end{enumerate}
      \end{description}
      %
      \begin{center}
            On donne $\ln 2 \simeq 0,7$ ; $\ln 3 \simeq 1,1$ ; $\ln 10 \simeq 2,3$ ;
            $\expo{} \simeq 2,7$ et $\expo{\frac{3}{2}} \simeq 4,5$.
      \end{center}
%%%%%%%%%%%%%%%%%%%%%%%%%%%%%%%%%%%%%%%%%%%%%%%%%%%%%%%%%%%%%%%%%%%%%%%%%%%%%%%%%%%%%%%%%%%%%%%%%%%%%%%%%%%%%%%%%%%%%%%%%%%%%%%%




























